The following finite arithmetic facts can be established without primality of \(N\). Write \(\alpha=1+\sqrt2\), \(\lambda=\alpha^2\) and \(N+F\sqrt2=\alpha^K\) with \(K\) an odd prime. The Pell identity gives \(N^2-2F^2=-1\), hence \(F\) is invertible modulo \(N\). Set \[
 s=F^{-1}\pmod N,\quad s^2=2,\quad h=3+2s.
\] Since \((1+s)^K=N+Fs=1\), \(h^K=1\). Also \(h-1=2(1+s)\) is a unit. Modulo each prime divisor of \(N\), the order of \(h\) is exactly \(K\). Thus \(1-h^j\) is a unit for \(1\leq j<K\), and \(1+h^j\) is a unit for every \(j\).

The symplectic matrix \[
 P=\begin{pmatrix}1&-1/(2s)\\s&1/2\end{pmatrix}
\] satisfies \(T=P\operatorname{diag}(h,h^{-1})P^{-1}\). Its lower-left entry is a unit, so its finite modular Weyl implementer has the same type of quadratic kernel as (LC.9) and preserves parity. Conjugating \(S_h^{-1}\) by this implementer constructs a phase-normalized \(U_N\) with \(U_N^K=I\). Here \(S_hu(x)=u(h^{-1}x)\). The unit conditions above show that every nonzero residue has an orbit of length \(K\), and that an orbit and its negative are distinct. On the odd space, there are \(n=(N-1)/(2K)\) cyclic orbits, each contributing every \(K\)th root once. Hence the multiplicity is exactly \(n\) for each root. A different global phase merely rotates all centres and changes none of the spacing conclusions.

Let \(W=A_{{\kappa_0},N}U_NA_{{\kappa_0},N}^*\) on the odd space. Restriction cannot increase (LC.20). For large \(N\) take \(\delta_N<\sin(\pi/K)\). Spectral stability for a normal matrix places each eigenvalue of \(V\) within \(\delta_N\) of a \(K\)th root. To verify multiplicities, let \(E=VW^*\) and use its principal Hermitian logarithm, \(E=e^{iH}\). Then \(\|H\|\leq2\arcsin(\delta_N/2)\), and the path \(e^{iuH}W\), \(0\leq u\leq1\), stays within \(\delta_N\) of \(W\). The disjoint slightly enlarged root disks have constant Riesz ranks along this path. Thus every disk contains exactly \(n\) eigenvalues. Their circle-coordinate halfwidth is bounded by \[
 a_N=\pi^{-1}\arcsin(\delta_N/2)=O(N^{-\infty}).              \tag{LC.22}
\] This is an operator-norm, individual-cluster statement, unlike the earlier normalized Hilbert--Schmidt comparison. It still makes no claim about individual eigenvectors inside a nearly degenerate cluster.

\subsection{Gaps, local counts, and robustness against unfolding}\label{locality:gaps-local-counts-and-robustness-against-unfolding}

Sort the \(D\) eigenangles cyclically, retaining multiplicities, and put \(S_l=D(\theta_{l+1}-\theta_l)\), with the cyclic final gap. Exactly \(D-K\) gaps are within clusters and at most \(2Da_N\). The other \(K\) gaps are in \([n-2Da_N,n+2Da_N]\). Matching these two classes to zeros and copies of \(n\) proves (LC.6). Since \(1/n=K/D\to0\), the reference measures in (LC.6) converge weakly to \(\delta_0\), and so do the actual gap measures. Their means are nevertheless exactly one at each \(N\): rare large gaps carry that mean.

Define the ordered close-pair count \[
 R_N(x)=\frac1D\#\{(l,j):l\ne j,\ D\operatorname{dist}_{\mathbb T}
                         (\theta_l,\theta_j)\leq x\}.
\] For fixed \(x>0\), eventually \(2Da_N<x\) and \(D(1/K-2a_N)>x\). Every within-cluster ordered pair is counted and no cross-cluster pair is counted. Therefore \[
 \boxed{R_N(x)=n-1\quad\hbox{eventually}.}                   \tag{LC.23}
\] For a fixed compactly supported continuously differentiable test \(\phi\), its periodized factorial statistic likewise obeys \[
 \mathcal R_N(\phi)=(n-1)\phi(0)
       +O\bigl((n-1)Da_N\|\phi'\|_\infty\bigr).             \tag{LC.24}
